\section{Discussion and open problems}\label{sec:discussion}

\subsection{What has been shown}

A passive system can delay a band of frequencies for free, but it cannot
shift the phase of the whole band by a constant without absorbing. The
least mean absorption is $\theta\Lambda$, with $\Lambda=(K(\kappa)-E(\kappa))/\kappa$
an explicit function of the band ratio, and the bound is attained. The proof
is a single sum rule, \eqref{eq:sumrule}, built from the test function $W$,
and every other statement of the paper is a consequence of it or of the
explicit extremal system: the lead bound, the worst-frequency and
geometric-mean bounds, the bound for anomalous dispersion in a medium, and
the stability under perturbations of the phase.

The ingredients are classical: the Herglotz representation, Kramers--Kronig
relations, a Cauchy integral, and an elliptic integral reduced by Landen's
transformation. What is new, as far as we can tell, is the inequality and
its sharpness, and the identification of the price of a constant phase with
a single number $\Lambda$.

\subsection{Open problems}

\begin{enumerate}
\item[(1)] \emph{Systems with zeros.} Determine the infimum of
\cref{prob:zeros}. Even its positivity is open. A first step would be the
best single zero, at an arbitrary point of $\Cp$ rather than on the
imaginary axis, and the best pair of zeros $\zeta,-\bar\zeta$.
\item[(2)] \emph{The plain mean for a lead.} \Cref{thm:lead} is sharp for the
weighted mean. Is the system $H_\star(-ab/z)$ of \cref{rem:leadmean}
optimal for the plain mean, with value $\abs\theta(E-\kappa'^2K)/\kappa$?
\item[(3)] \emph{The worst frequency.} \Cref{cor:worst} bounds $\max\alpha$ by
the mean. The sharp bound on $\max\alpha$ is an extremal problem of
Chebyshev type: minimise $\sup_{[a,b]}\Im h$ among Herglotz functions with
$\Re h=\theta$ on the band. The extremal $H_\star$ has
$\max\alpha_\star/\theta\Lambda=1.273$ for an octave, so the sharp constant
lies between $\Lambda$ and $1.273\Lambda$ there.
\item[(4)] \emph{Several bands.} For a phase that is constant on each of
several disjoint bands, the test function is built from a hyperelliptic
square root, and the constant becomes a period of a hyperelliptic
differential. We expect the sum rule to carry over with the same proof.
\item[(5)] \emph{Reflection and matching.} For reflection coefficients,
the Fano bound \cite{fano} and the present bound constrain different parts
of the same function. A joint bound on the modulus and the argument of a
reflection coefficient on a band would combine the two.
\item[(6)] \emph{Formalisation.} The finite arithmetic of the paper is
checked in Lean (\cref{app:verification}). The analytic part, the sum rule
and the evaluation of $\Lambda$, would require a formalisation of boundary
values of Cauchy integrals, of the M.~Riesz theorem and of the complete
elliptic integrals with Landen's transformation. We have not done this, and
the analytic theorems of the paper are not machine-checked.
\end{enumerate}
